%% ma: L12-B08-0009
%% lop: 12
%% chuong: 2
%% bai: 8
%% dang: 12.8.6
%% loai: TLN
%% muc: VD
%% dap_an: "3"
%% nguon: "(NBV)-ĐỀ SỐ 4-MỖI NGÀY 1 ĐỀ THI 2025.pdf (D04, MỖI NGÀY 1 ĐỀ - Đề số 4), Phần III Câu 6"
%% kien_thuc: "Bài 7-8 (hệ trục tọa độ, độ dài vectơ), Bài 16 (góc giữa đường thẳng và mặt phẳng); tỉ số lượng giác (Lớp 9)"
%% trang_thai: chinh-thuc
%% ghi_chu: "Hình gốc chỉ minh họa bố trí hai bức tường nên không vẽ lại (đề đã mô tả đủ). Đáp án gốc 3 đúng (sympy: 3,026)."
\begin{ex}
Một sợi dây được căng giữa hai bức tường (vuông góc với nhau và cùng vuông góc với mặt đất). Đầu $A$ của sợi dây nằm trên bức tường thứ nhất, cách bức tường thứ hai $3\ \mathrm m$ và cách mặt đất $1{,}2\ \mathrm m$. Đầu $B$ cao hơn đầu $A$ và nằm trên bức tường thứ hai, cách bức tường thứ nhất $1\ \mathrm m$. Biết rằng góc giữa dây căng và mặt đất không vượt quá $30^\circ$, hỏi khoảng cách từ $B$ đến mặt đất tối đa bao nhiêu mét? Kết quả làm tròn đến hàng đơn vị.
\shortans{3}
\loigiai{Chọn hệ trục $Oxyz$ có $O$ là chân giao tuyến của hai bức tường, $Ox$ dọc theo tường thứ nhất, $Oy$ dọc theo tường thứ hai, $Oz$ thẳng đứng; mặt đất là $(Oxy)$. Khi đó $A(3;0;1{,}2)$ và $B(0;1;z_B)$ với $z_B>1{,}2$.
Gọi $h=z_B-1{,}2>0$. Hình chiếu của $AB$ lên mặt đất có độ dài $\sqrt{3^2+1^2}=\sqrt{10}$, nên góc $\alpha$ giữa dây và mặt đất thoả $\tan\alpha=\dfrac h{\sqrt{10}}\le\tan30^\circ=\dfrac1{\sqrt3}$, suy ra $h\le\sqrt{\frac{10}3}\approx1{,}83$.
Vậy $z_B\le1{,}2+1{,}83\approx3{,}03$, tức khoảng cách tối đa từ $B$ đến mặt đất xấp xỉ $3\ \mathrm m$.}
\end{ex}
